\section{Circom}
\label{sec:circom}

La Sección~\ref{sec:computacion_prueba} presentó el enfoque basado en circuitos aritméticos y
restricciones, en el cual la computación que se desea probar se expresa como un conjunto de
ecuaciones sobre un cuerpo finito. Circom constituye una de las herramientas más difundidas
dentro de ese enfoque y permite observar en concreto cómo una descripción de la computación se
transforma en un sistema de restricciones, cómo se obtiene el testigo correspondiente y cómo ambos
elementos se utilizan para generar una prueba. La presente sección describe el lenguaje y su
compilador, el proceso que conduce desde un circuito hasta una prueba verificable en Ethereum y
las limitaciones que este modelo de desarrollo impone, en particular cuando se aplica a
verificaciones propias de las redes blockchain.

% ---------------------------------------------------------------------------
\subsection{Origen y propósito}
\label{sec:circom_origen}

Circom es un lenguaje para la descripción de circuitos aritméticos, acompañado de un compilador
que traduce dichas descripciones a un sistema de restricciones de rango uno (R1CS). Fue
desarrollado en el ecosistema del proyecto iden3, y su diseño y semántica fueron descritos
académicamente por Bellés-Muñoz et al.~\cite{belles2023circom}. La versión actual del
compilador, conocida como Circom~2, está implementada en Rust y reemplazó a una primera
implementación en JavaScript que se mantiene congelada~\cite{circomdocs}. El ecosistema se
completa con circomlib, una biblioteca de plantillas reutilizables que implementan primitivas
criptográficas como las funciones de hash Poseidon, MiMC y Pedersen, firmas EdDSA, comparadores
y árboles de Merkle dispersos~\cite{circomlib}, y con snarkjs, una implementación en JavaScript
y WebAssembly de sistemas de prueba que consume los artefactos producidos por el
compilador~\cite{snarkjs}.

Circom es un \textbf{lenguaje de dominio específico} (\textit{domain-specific language}, DSL), es
decir, un lenguaje cuyas construcciones corresponden directamente a los elementos de un dominio
de problemas concreto, en lugar de ofrecer las abstracciones generales de un lenguaje de
programación convencional. En este caso, los elementos del dominio son las señales, que
representan valores de $\mathbb{F}_p$, y las restricciones que dichas señales deben satisfacer.
Esta orientación determina la principal diferencia entre escribir un programa convencional y
describir un circuito en Circom. En un programa convencional, las instrucciones se ejecutan
secuencialmente y su efecto queda determinado por el resultado de esa ejecución. En un circuito
de Circom, en cambio, la descripción cumple dos funciones distintas. Por un lado, indica cómo
calcular los valores de las señales a partir de las entradas, cálculo que solo realiza el
probador. Por otro, establece las ecuaciones que esos valores deben satisfacer, y únicamente
estas ecuaciones forman parte de lo que el verificador comprueba. La separación entre ambas
funciones, que Circom expone de forma explícita al desarrollador, es el aspecto central del
lenguaje y la fuente de sus principales riesgos, analizados en la
Sección~\ref{sec:circom_limitaciones}.

Circom no constituye por sí mismo un sistema de prueba. Su compilador produce el sistema de
restricciones y un programa capaz de calcular el testigo, pero la configuración inicial, la
generación de la prueba y su verificación corresponden a herramientas externas, como snarkjs, que
implementan los sistemas descritos en la Sección~\ref{sec:snarks_starks}.

% ---------------------------------------------------------------------------
\subsection{Construcción de circuitos}
\label{sec:circom_construccion}

La unidad básica de descripción en Circom es la \textbf{plantilla} (\textit{template}), que
define un subcircuito parametrizable con un conjunto de señales y las restricciones que las
vinculan. Una plantilla puede instanciarse como \textbf{componente} dentro de otra plantilla,
lo que permite construir circuitos complejos a partir de piezas más simples, del mismo modo en
que circomlib ofrece plantillas para funciones de hash o comparaciones que se integran en
circuitos mayores. Todo circuito tiene un componente principal, denominado \lstinline|main|, que
define la interfaz del circuito completo~\cite{circomdocs}.

Las \textbf{señales} son los valores sobre los que operan las restricciones y se clasifican en
señales de entrada (\lstinline|signal input|), de salida (\lstinline|signal output|) e
intermedias (\lstinline|signal|). Su visibilidad se relaciona directamente con las nociones de
enunciado y testigo de la Sección~\ref{sec:nocion_zk}. Las señales de salida del componente
principal son siempre públicas, las señales de entrada son privadas salvo que se declaren
explícitamente como públicas mediante la sintaxis \lstinline[breaklines=true]|component main {public [...]}|, y
las señales intermedias son siempre privadas~\cite{circomdocs}. Las señales públicas forman, por
lo tanto, el enunciado que conoce el verificador, mientras que los valores de todas las señales
conforman, conceptualmente, el testigo. Circom admite además variables declaradas con
\lstinline|var| que, a diferencia de las señales, no forman parte del testigo ni generan
restricciones por sí mismas. Una variable puede contener valores conocidos durante la
compilación, como índices o contadores, pero también valores que dependen de señales y que, por
lo tanto, solo se conocen durante el cálculo del testigo, así como expresiones aritméticas que el
compilador emplea para construir restricciones~\cite{circomdocs}.

La relación entre señales se establece mediante tres operadores, cuya diferencia es esencial
para comprender el lenguaje~\cite{circomdocs}. El operador \lstinline|===| impone una
restricción entre dos expresiones sin asignar ningún valor. El operador \lstinline|<--| asigna un
valor a una señal sin generar ninguna restricción, de modo que ese valor solo interviene en el
cálculo del testigo. El operador \lstinline|<==| combina ambas acciones, ya que asigna el valor
y genera simultáneamente la restricción de igualdad correspondiente. Los operadores
\lstinline|==>| y \lstinline|-->| son las variantes de \lstinline|<==| y \lstinline|<--| con la
asignación escrita en sentido inverso. El compilador exige que toda restricción sea de grado a
lo sumo dos y pueda llevarse a la forma $A \cdot B + C = 0$, donde $A$, $B$ y $C$ son
combinaciones lineales de señales, y rechaza las restricciones que no admiten esa
forma~\cite{circomdocs}. Esta forma coincide con la de una fila de un sistema R1CS, en el que
cada restricción relaciona el producto de dos combinaciones lineales con una tercera
(Sección~\ref{sec:computacion_prueba}). Las restricciones lineales se representan en este
formato tomando una constante como uno de los factores, mientras que no toda expresión de grado
dos admite una única fila. La restricción $e = a \cdot b + c \cdot d$, por ejemplo, es rechazada
por el compilador y debe descomponerse mediante una señal intermedia en dos restricciones. Cada
restricción de Circom tiene, por lo tanto, la forma de una fila de dicho sistema, aunque el
compilador puede simplificar y transformar el conjunto de restricciones antes de producirlo.

La Figura~\ref{fig:circom_cubico} muestra un ejemplo mínimo que retoma la computación utilizada
en la Sección~\ref{sec:computacion_prueba}. El circuito permite demostrar el conocimiento de un
valor privado $w$ tal que $w^{3} + w + 5 = y$, siendo $y$ una salida pública.

\begin{figure}[htbp]
	\centering
	\begin{minipage}{0.8\textwidth}
		\begin{lstlisting}[language=Circom]
pragma circom 2.1.6;

template Cubic() {
    signal input w;     // entrada privada
    signal output y;    // salida publica
    signal v1;          // senales intermedias
    signal v2;

    v1 <== w * w;
    v2 <== v1 * w;
    y  <== v2 + w + 5;
}

component main = Cubic();
\end{lstlisting}
	\end{minipage}
	\caption{Circuito Circom que expresa la relación $w^{3} + w + 5 = y$.}
	\label{fig:circom_cubico}
\end{figure}

Cada uno de los operadores \lstinline|<==| del ejemplo expresa una restricción, de modo que el
código describe el sistema
\begin{equation}
	\label{eq:circom_cubico}
	w \cdot w - v_1 = 0, \qquad v_1 \cdot w - v_2 = 0, \qquad (v_2 + w + 5) \cdot 1 - y = 0,
\end{equation}
cuyas dos primeras restricciones son no lineales y cuya tercera es lineal. Todas las operaciones
se realizan en $\mathbb{F}_p$, donde $p$ es, por defecto, el orden primo de 254 bits del grupo
asociado a la curva BN254, la misma curva para la que Ethereum dispone de contratos
precompilados~\cite{circomdocs, eip197}. Al evaluar el circuito para la entrada $w = 3$ se
obtienen los valores $(y, w, v_1, v_2) = (35, 3, 9, 27)$, que corresponden a la asignación
presentada en la Sección~\ref{sec:computacion_prueba}. Solo $y$ es una señal pública, por lo
que su valor forma parte de las entradas públicas del sistema de prueba, mientras que $w$ queda
excluida de ellas. Esta exclusión determina únicamente qué valores recibe el verificador. Que la
prueba no revele información adicional sobre $w$ más allá de la que pueda deducirse del valor
público $y$, depende de que el sistema de prueba empleado satisfaga la propiedad de conocimiento 
cero descrita en la Sección~\ref{sec:propiedades_zk}.

El sistema~\eqref{eq:circom_cubico} reúne las restricciones expresadas en el código, que no
coinciden necesariamente con las que conserva el sistema compilado. El compilador ofrece tres
niveles de simplificación. La opción \lstinline|--O0| no aplica ninguna, \lstinline|--O1| solo
simplifica igualdades entre señales y entre señales y constantes, y \lstinline|--O2| aplica la
simplificación completa, que puede eliminar restricciones lineales sustituyendo una señal por la
expresión que la define~\cite{circomdocs}. El nivel aplicado por defecto depende de la versión
del compilador, ya que fue \lstinline|--O2| hasta la versión 2.1.9 y pasó a ser \lstinline|--O1|
a partir de la versión 2.2.0~\cite{circomreleases}. La simplificación completa despeja, de cada
restricción lineal que contenga una señal privada, una de esas señales y la sustituye en las
demás restricciones, sin eliminar nunca las señales públicas~\cite{circomdocs}. En una
compilación del ejemplo realizada para el presente trabajo con la versión 2.2.3, las opciones
\lstinline|--O0| y \lstinline|--O1| conservaron las tres restricciones, mientras que
\lstinline|--O2| despejó $v_2 = y - w - 5$ de la restricción lineal y lo sustituyó en la
segunda, de modo que el sistema resultante constó de las dos restricciones $w \cdot w = v_1$ y
$v_1 \cdot w = y - w - 5$, en las que $v_2$ ya no interviene. La elección de la señal que se
despeja corresponde al algoritmo del compilador, por lo que el sistema resultante es una de
varias formas algebraicamente equivalentes de expresar la misma relación.

El ejemplo permite precisar qué significa que un circuito esté \textbf{correctamente
restringido}. Un circuito lo está cuando las asignaciones que satisfacen su sistema de
restricciones corresponden exactamente a los casos en que se cumple la relación que el
desarrollador pretende demostrar. Esta condición no exige que todas las señales auxiliares
queden determinadas unívocamente. En el circuito de la Figura~\ref{fig:circom_cubico}, una vez
fijado $w$, las restricciones determinan $v_1$, $v_2$ e $y$, de modo que ningún probador puede
presentar un valor de $y$ distinto de $w^{3} + w + 5$. En la plantilla \texttt{IsZero} de la
Sección~\ref{sec:circom_limitaciones}, en cambio, cuando la entrada es nula las restricciones
fijan la salida en $1$ pero admiten cualquier valor de la señal auxiliar \texttt{inv}, libertad
que no permite aceptar ninguna afirmación incorrecta. Un circuito \textit{insuficientemente
restringido} admite, por el contrario, asignaciones que conducen a aceptar afirmaciones que la
especificación rechaza, mientras que un circuito \textit{sobrerrestringido} rechaza asignaciones
que la especificación considera legítimas. Estos defectos deben distinguirse de las propiedades
del sistema de prueba descritas en la Sección~\ref{sec:propiedades_zk}. El sistema criptográfico
sigue siendo sólido y completo respecto de la relación que efectivamente se compiló, pero un
circuito insuficientemente restringido compromete la solidez respecto de la especificación
prevista, y uno sobrerrestringido impide demostrar afirmaciones que, según dicha especificación,
deberían poder demostrarse.

% ---------------------------------------------------------------------------
\subsection{Compilación, cálculo del testigo y generación de la prueba}
\label{sec:circom_flujo}

El proceso que conduce desde un circuito de Circom hasta una prueba verificable concreta las
etapas generales descritas en la Sección~\ref{sec:computacion_prueba}, y se distribuye entre el
compilador de Circom y un sistema de prueba externo. La Figura~\ref{fig:circom_flujo} representa
este proceso para el caso de Groth16 implementado en snarkjs.

En la primera etapa, el compilador procesa el circuito y genera dos artefactos independientes.
El primero es el sistema de restricciones, almacenado en un archivo binario con extensión
\texttt{.r1cs}, que define formalmente la relación que la prueba certificará. El segundo es un
programa capaz de calcular el testigo, que el compilador genera en WebAssembly o en C++, siendo
esta última opción más eficiente para circuitos de gran tamaño~\cite{circomdocs}. Durante la
compilación, el compilador comprueba además que todas las restricciones sean de grado a lo sumo
dos y admitan la forma descrita, y aplica las simplificaciones correspondientes al nivel de
optimización elegido. Desde la versión 2.0.8 incorpora también una advertencia para los usos de
\lstinline|<--| en los que la expresión asignada es cuadrática y podría emplearse
\lstinline|<==|~\cite{circomreleases}. En la versión 2.2.3 utilizada en las compilaciones de
este trabajo, dicha advertencia se emitió al activar la opción \lstinline|--inspect|, que realiza
comprobaciones adicionales destinadas a identificar señales potencialmente insuficientemente
restringidas, como aquellas que no aparecen en ninguna restricción~\cite{circomdocs}.

\begin{figure}[htbp]
	\centering
	\begin{tikzpicture}[
		herramienta/.style = {draw, rounded corners, align=center, minimum height=1.1cm,
			text width=2.2cm, font=\footnotesize},
		artefacto/.style = {draw, dashed, align=center, minimum height=0.9cm,
			text width=2.2cm, font=\footnotesize},
		flecha/.style = {-{Stealth[length=2.5mm]}, thick},
		dato/.style = {font=\footnotesize, align=center}
	]
		\node[artefacto] (fuente) at (0, 0) {Circuito\\\texttt{.circom}};
		\node[herramienta] (compilador) at (2.65, 0) {Compilador\\circom};
		\node[artefacto] (r1cs) at (5.3, 1.5) {Restricciones\\\texttt{.r1cs}};
		\node[artefacto] (wasm) at (5.3, -1.5) {Calculador\\del testigo};
		\node[herramienta] (setup) at (7.95, 1.5) {Configuración\\inicial};
		\node[herramienta] (calc) at (7.95, -1.5) {Cálculo del\\testigo};
		\node[artefacto] (entradas) at (7.95, -3.5) {Entradas\\\texttt{input.json}}; 
		\node[artefacto] (zkey) at (10.6, 1.5) {Claves\\\texttt{.zkey}};
		\node[artefacto] (wtns) at (10.6, -1.5) {Testigo\\\texttt{.wtns}};
		\node[herramienta] (prove) at (13.25, -1.5) {Generación\\de la prueba};
		\node[herramienta] (verify) at (13.25, -3.9) {Verificación};

		\draw[flecha] (fuente) -- (compilador);
		\draw[flecha] (compilador.north) |- (r1cs.west);
		\draw[flecha] (compilador.south) |- (wasm.west);
		\draw[flecha] (r1cs) -- (setup);
		\draw[flecha] (setup) -- (zkey);
		\draw[flecha] (wasm) -- (calc);
		\draw[flecha] (entradas) -- (calc);
		\draw[flecha] (calc) -- (wtns);
		\draw[flecha] (wtns) -- (prove);
		\draw[flecha] (zkey.east) -| node[right, dato, pos=0.75] {clave de\\prueba} (prove.north);
		\draw[flecha] (prove) -- node[left, dato] {$\pi$ y señales\\públicas} (verify);
		\draw[flecha] (zkey.north) -- ++(0, 0.5) -| node[above, dato, pos=0.25]
			{clave de verificación} ++(4.3, -2.0) |- (verify.east);
	\end{tikzpicture}
	\caption{Proceso de compilación, cálculo del testigo, generación y verificación de una
	prueba a partir de un circuito Circom.}
	\label{fig:circom_flujo}
\end{figure}

En la segunda etapa, el probador ejecuta el calculador del testigo sobre entradas concretas,
expresadas en un archivo en formato JSON, y obtiene un archivo \texttt{.wtns} con la asignación correspondiente al sistema R1CS
compilado~\cite{circomdocs}. Este archivo contiene los valores de las señales que el sistema
conserva, precedidos por la constante $1$ que el formato emplea para representar términos
independientes, por lo que su tamaño y su orden dependen de la compilación. En el ejemplo de la
Figura~\ref{fig:circom_cubico}, el testigo obtenido con \lstinline|--O1| contiene los cinco
valores $(1, 35, 3, 9, 27)$, mientras que con \lstinline|--O2| la señal $v_2$ no forma parte del
sistema y el testigo se reduce a cuatro valores. Este cálculo
ejecuta las asignaciones descritas mediante los operadores \lstinline|<--| y \lstinline|<==|,
pero no constituye una garantía criptográfica, ya que un probador puede producir un archivo de
testigo por cualquier otro medio. Lo que determina si una prueba será aceptada no es el programa
que calculó el testigo, sino el hecho de que dicho testigo satisfaga el sistema de restricciones.

La tercera etapa corresponde a la configuración inicial del sistema de prueba. En el caso de
Groth16, cuyos fundamentos se presentaron en la Sección~\ref{sec:snarks_starks}, la
configuración se divide en dos fases. La primera, denominada \textit{powers of tau}, es
independiente del circuito y puede reutilizarse para cualquier circuito hasta un tamaño máximo.
La segunda es específica del circuito y produce un archivo \texttt{.zkey} que contiene la clave
de prueba, a partir del cual se exporta la clave de verificación~\cite{circomdocs}. Ambas fases
admiten contribuciones de múltiples participantes, en el sentido de las ceremonias descritas en
la Sección~\ref{sec:snarks_starks}~\cite{bowe2017scalable}. snarkjs implementa también PLONK y
FFLONK, que solo requieren la fase universal~\cite{snarkjs}.

En la cuarta etapa, el probador combina la clave de prueba y el testigo para generar la prueba
$\pi$, junto con un archivo que contiene los valores de las señales públicas. Finalmente, la
verificación puede realizarse con snarkjs o mediante un contrato inteligente en Solidity que la
propia herramienta genera a partir de la clave de verificación~\cite{snarkjs}. Dado que, por
defecto, Circom trabaja sobre el cuerpo asociado a la curva BN254, el contrato generado puede
apoyarse en los contratos precompilados de Ethereum para dicha curva~\cite{eip196, eip197}, lo
que hace posible la verificación \textit{on-chain} descrita en la
Sección~\ref{sec:computacion_prueba}.

Este proceso permite delimitar las responsabilidades de cada componente. El compilador define
la relación que se probará, el calculador del testigo obtiene una asignación que la satisface, el
sistema de prueba genera a partir de ella una prueba sucinta y el verificador comprueba esa
prueba. En el caso de Groth16, la decisión del verificador se basa en la clave de verificación,
la prueba $\pi$ y los valores de las entradas públicas. El verificador no recibe el archivo
\texttt{.r1cs} ni ejecuta el programa que calculó el testigo, pero la clave de verificación se
deriva del sistema de restricciones compilado durante la configuración inicial, por lo que la
relación que la prueba certifica es exactamente la que dicho sistema expresa. El programa de
generación del testigo no aporta, por sí mismo, ninguna garantía criptográfica, y cualquier
propiedad que no haya quedado expresada como restricción carece de respaldo, con independencia
de lo que haga ese programa.

% ---------------------------------------------------------------------------
\subsection{Limitaciones y consideraciones de seguridad}
\label{sec:circom_limitaciones}

La separación entre cálculo y restricción descrita en la Sección~\ref{sec:circom_origen} es la
principal fuente de errores en el desarrollo con Circom. El operador \lstinline|<--| resulta
necesario cuando el cálculo del valor de una señal no puede expresarse directamente mediante una
restricción cuadrática,
pero el valor así asignado solo queda respaldado por la prueba si el desarrollador añade
restricciones que lo vinculen con las demás señales. La plantilla \texttt{IsZero} de circomlib,
reproducida en la Figura~\ref{fig:circom_iszero}, ilustra este patrón~\cite{circomlib}.

\begin{figure}[htbp]
	\centering
	\begin{minipage}{0.8\textwidth}
		\begin{lstlisting}[language=Circom]
template IsZero() {
    signal input in;
    signal output out;
    signal inv;

    inv <-- in != 0 ? 1 / in : 0; // solo calculo
    out <== -in * inv + 1;         // asigna y restringe
    in * out === 0;                // solo restringe
}
\end{lstlisting}
	\end{minipage}
	\caption{Plantilla \texttt{IsZero} de circomlib, que combina un cálculo auxiliar con
	restricciones explícitas.}
	\label{fig:circom_iszero}
\end{figure}

La plantilla devuelve $1$ si la entrada es nula y $0$ en caso contrario. Para una entrada
distinta de cero, la relación $\texttt{in} \cdot \texttt{inv} = 1$ es cuadrática, pero la
expresión que calcula \texttt{inv} es condicional y emplea una división, por lo que no puede
asignarse mediante \lstinline|<==|. Por ese motivo, la asignación
\lstinline[breaklines=true]|inv <-- in != 0 ? 1 / in : 0| solo calcula un valor candidato, que se
incorpora al testigo, mientras que las dos restricciones establecen las relaciones que toda
asignación aceptada debe cumplir. Si la entrada es distinta de cero, la segunda restricción
obliga a que la salida sea nula, y la primera obliga entonces a que \texttt{inv} sea el inverso
de la entrada. Si la entrada es nula, la primera restricción obliga a que la salida sea $1$, sin
determinar un único valor de \texttt{inv}. La garantía no consiste, por lo tanto, en que
cualquier valor de \texttt{inv} satisfaga las restricciones, sino en que toda asignación que las
satisface produce la salida correcta. El valor calculado sin restricción actúa como una
sugerencia que las restricciones verifican. Si, en cambio, un circuito asignara su salida mediante
\lstinline|y <-- w * w;| sin ninguna restricción adicional, el compilador no generaría ninguna
ecuación que vincule $y$ con $w$, y un probador deshonesto podría generar una prueba válida para
cualquier valor de $y$. En este caso concreto, las comprobaciones descritas en la
Sección~\ref{sec:circom_flujo} emiten advertencias, tanto por el uso de \lstinline|<--| con una
expresión cuadrática como por la presencia de señales que no aparecen en ninguna
restricción~\cite{circomdocs}. Estas comprobaciones identifican patrones conocidos, pero no
detectan todas las vulnerabilidades posibles ni sustituyen una revisión de seguridad del
circuito, ya que determinar en general si un circuito está correctamente restringido requiere
análisis específicos como los que se describen a continuación.

Los circuitos insuficientemente restringidos constituyen una de las categorías de
vulnerabilidades más estudiadas en sistemas de prueba sucintos~\cite{chaliasos2024sok}, y han
motivado el desarrollo de herramientas de análisis específicas. Pailoor et al.\ propusieron un
método automático para detectar circuitos insuficientemente restringidos~\cite{pailoor2023automated},
y Wen et al.\ desarrollaron un marco de análisis estático para circuitos de Circom que, evaluado
sobre más de 258 circuitos de proyectos de código abierto, detectó vulnerabilidades previamente
desconocidas~\cite{wen2024practical}. Conviene subrayar que este riesgo no es exclusivo de Circom,
sino propio de todo enfoque en el que el desarrollador define explícitamente las restricciones.
Circom lo hace visible porque expone de manera directa la diferencia entre asignar un valor y
restringirlo.

Una segunda limitación, también común al enfoque basado en circuitos, es el costo de las
operaciones que no son nativas del cuerpo finito, señalado en la
Sección~\ref{sec:computacion_prueba}. La plantilla \texttt{Num2Bits} de circomlib, que
descompone un valor en $n$ bits, introduce una restricción de la forma $b_i (b_i - 1) = 0$ para
cada bit y una restricción adicional que recompone el valor original~\cite{circomlib}. Las
funciones de hash que operan sobre palabras binarias, como SHA-256 y Keccak-256, deben
construirse sobre descomposiciones de este tipo, lo que multiplica el número de restricciones. El
formato R1CS, además, no incorpora de forma nativa los argumentos de búsqueda mencionados en la
Sección~\ref{sec:computacion_prueba}, que se retoman en la Sección~\ref{sec:noir}.

Estas limitaciones adquieren especial relevancia en las verificaciones propias de los puentes
entre cadenas, en las que conviene distinguir tres afirmaciones diferentes. La primera es que una
transacción fue incluida en un bloque, lo que se demuestra mediante una prueba de inclusión
respecto de la raíz de transacciones del encabezado y exige recalcular hashes Keccak-256 sobre
los nodos del árbol de Merkle-Patricia y sobre el propio encabezado. La segunda es que dicho
bloque es auténtico y está finalizado según el consenso de la cadena de origen. La tercera es que
la operación de bloqueo se ejecutó con éxito, ya que una transacción incluida en un bloque puede
haber revertido su ejecución. La prueba de inclusión no establece por sí misma esta última
propiedad, cuya demostración puede requerir evidencia adicional, como el recibo de la
transacción, que registra su estado de ejecución y los eventos emitidos, junto con su inclusión
respecto de la raíz de recibos del encabezado~\cite{wood2014ethereum, eip658}. La segunda
afirmación, por su parte, no se reduce a la verificación de una firma. El protocolo de cliente
ligero de Ethereum exige, además de verificar la firma agregada del comité de sincronización,
comprobar el número de miembros participantes, el dominio de firma correspondiente a la versión
del protocolo, la pertenencia del encabezado finalizado al estado certificado mediante una rama
de Merkle y la continuidad del propio comité entre periodos~\cite{ethlightclient}. A partir de
la actualización Capella, el protocolo incorpora además la vinculación entre el encabezado de
consenso y el encabezado de ejecución que contiene las raíces
mencionadas~\cite{ethlightclientcapella}. Entre estas
comprobaciones, la verificación de la firma es la más costosa de expresar como circuito. Esta
firma emplea la curva BLS12-381, cuyo cuerpo base
tiene 381 bits y excede, por lo tanto, al cuerpo de 254 bits sobre el que opera el circuito, de
manera que cada elemento debe representarse mediante varios valores del cuerpo y cada operación
aritmética debe emularse. La implementación de referencia circom-pairing, desarrollada en el
marco de un grupo de trabajo de 0xPARC, documenta que la verificación de una única firma BLS
requiere aproximadamente 19,2 millones de restricciones no lineales, aproximadamente 3,2 horas de
compilación y una clave de prueba de 12~GB en un servidor de 32 núcleos y 256~GB de
memoria~\cite{circompairing}. Estas cifras, obtenidas sobre un hardware concreto, no constituyen
una medición del presente trabajo, pero ilustran el orden de magnitud del costo que implica
expresar verificaciones de consenso como circuitos aritméticos. La forma en que el puente
desarrollado en este trabajo aborda cada una de las tres afirmaciones se expone en el
Capítulo~\ref{chp:zkbridge}.

Circom ofrece, en definitiva, un control preciso sobre las restricciones generadas, lo que
permite optimizar manualmente los circuitos y conocer con exactitud la relación que se probará.
Ese control tiene como contrapartida una mayor dificultad de desarrollo, ya que el programador
debe razonar simultáneamente sobre el cálculo del testigo y sobre las ecuaciones que lo
restringen, y asume la responsabilidad de que ambos coincidan. Noir, analizado en la
Sección~\ref{sec:noir}, conserva el enfoque basado en circuitos, pero eleva el nivel de
abstracción, de modo que el desarrollador describe la computación con construcciones más
cercanas a las de un lenguaje de programación convencional y el compilador asume la generación de
las restricciones correspondientes.
